\documentclass[11pt,a4paper]{article}

\usepackage[utf8]{inputenc}
\usepackage[T1]{fontenc}
\usepackage[margin=2.5cm]{geometry}
\usepackage{amsmath,amssymb}
\usepackage{booktabs}
\usepackage{multirow}
\usepackage{graphicx}
\usepackage{hyperref}
\usepackage{array}
\usepackage{xcolor}
\usepackage{caption}
\usepackage{microtype}
\usepackage{parskip}
\usepackage{float}

\hypersetup{
  colorlinks=true,
  linkcolor=blue!70!black,
  urlcolor=blue
}

\title{\textbf{Local Search Algorithms for the Linear Ordering Problem:\\
An Empirical Comparison of Neighborhoods, Pivoting Rules, and Initialization Strategies}}
\author{Daniele Cioffi\\
\small Heuristic Optimization --- Exercise 1}
\date{April 2026}

\begin{document}

\maketitle

%─────────────────────────────────────────────────────────────────────────────
\begin{abstract}
We implement and compare 14 local search configurations for the
\emph{Linear Ordering Problem} (LOP), combining three neighborhood
structures (Transpose, Exchange, Insert), two pivoting rules
(First Improvement, Best Improvement), two initialization strategies
(Random, Chenery--Watanabe), and two Variable Neighborhood Descent
(VND) variants.  Experiments are conducted on the full XLOLIB benchmark
of 78 instances (sizes $n=150$ and $n=250$).  Results are evaluated via
the Relative Percentage Deviation (apd) from the best-known objective
value.  Statistical significance is assessed with paired $t$-tests and
Wilcoxon signed-rank tests.  The Insert neighborhood yields dramatically
lower apd (mean~$\approx 5\%$) compared to Transpose and Exchange
(mean~$>15\%$), a difference that is statistically significant across
all comparisons.  The Chenery--Watanabe initialization is significantly
beneficial for the weaker neighborhoods but confers no clear advantage
for Insert.  Both VND variants match the quality of
\textsc{First}+\textsc{Transpose}+\textsc{Chenery} but never surpass
Insert-based search.
\end{abstract}

%─────────────────────────────────────────────────────────────────────────────
\section{Introduction}

The Linear Ordering Problem (LOP) is a well-known combinatorial
optimization problem with applications in economic input--output
analysis, triangulation of matrices, archaeological sequencing, and
machine scheduling.
It is NP-hard in the general case, making exact methods impractical for
large instances, and it has served as a canonical benchmark for
evaluating local search and metaheuristic techniques.

This report investigates how three design choices interact in
deterministic local search for the LOP:
\begin{enumerate}
    \item \textbf{Neighborhood structure} --- how the set of neighbors of
          a solution is defined;
    \item \textbf{Pivoting rule} --- whether to accept the \emph{first}
          improving neighbor or to scan the full neighborhood and accept
          the \emph{best};
    \item \textbf{Initial solution} --- whether to start from a random
          permutation or from the Chenery--Watanabe (CW) greedy
          construction heuristic.
\end{enumerate}
In addition, two Variable Neighborhood Descent (VND) strategies that
cycle through all three neighborhoods are evaluated.

All algorithms are implemented in C and run on the XLOLIB benchmark set.
The remainder of this report is organized as follows.
Section~\ref{sec:problem} formalizes the LOP.
Section~\ref{sec:methods} describes the implementation.
Section~\ref{sec:results} presents experimental results and statistical
analysis.
Section~\ref{sec:conclusion} concludes.

%─────────────────────────────────────────────────────────────────────────────
\section{Problem Description}
\label{sec:problem}

Given an $n\times n$ matrix $C=(c_{ij})$ of non-negative real weights,
the LOP asks for a permutation $\pi$ of $\{1,\ldots,n\}$ that maximizes

\begin{equation}
  f(\pi) \;=\; \sum_{i=1}^{n}\sum_{j=i+1}^{n} c_{\pi_i,\pi_j}.
  \label{eq:lop}
\end{equation}

That is, we seek an ordering of the rows (and columns) of $C$ such that
the sum of entries in the upper triangle of the permuted matrix is
maximized.  Equivalently, the LOP corresponds to finding an acyclic
tournament that maximizes the weight of forward arcs.

\paragraph{XLOLIB benchmark.}
The experiments use the XLOLIB instance set, which comprises instances
derived from real-world economic input--output tables.  The full benchmark comprises
78~instances: 39 with $n=150$ and 39 with $n=250$.  Best-known
objective values are provided for each instance and used to compute the
Relative Percentage Deviation (apd):

\begin{equation}
  \mathrm{apd} \;=\; 100 \times \frac{f^*-f(\hat\pi)}{f^*},
  \label{eq:apd}
\end{equation}

where $f^*$ is the best-known value and $f(\hat\pi)$ is the value found
by the algorithm.  A lower apd indicates better solution quality.

%─────────────────────────────────────────────────────────────────────────────
\section{Materials and Methods}
\label{sec:methods}

\subsection{Implementation}

The solver is implemented in C (standard C99) and organized into five
modules: argument parsing and main loop (\texttt{main.c}), instance
loading (\texttt{instance.c}), all search algorithms
(\texttt{optimization.c}), a custom linear congruential random number
generator (\texttt{utilities.c}), and wall-clock timing
(\texttt{timer.c}).

The RNG seed is set deterministically as the sum of all entries in the
cost matrix, ensuring reproducibility without a fixed external seed.

\subsection{Neighborhood Structures}

Three standard LOP neighborhoods are implemented.

\paragraph{Transpose ($\mathcal{N}_T$).}
Swap two \emph{adjacent} elements at positions $i$ and $i+1$.  The
neighborhood size is $n-1$.  The incremental cost change for transposing
elements $\pi_i$ and $\pi_{i+1}$ is
$\Delta = c_{\pi_{i+1},\pi_i} - c_{\pi_i,\pi_{i+1}}$,
computed in $O(1)$.

\paragraph{Exchange ($\mathcal{N}_E$).}
Swap any two elements at positions $i$ and $j$ ($i < j$).  The
neighborhood size is $\binom{n}{2}$.  The incremental cost change
(\texttt{deltaExchange}) is computed in $O(n)$.

\paragraph{Insert ($\mathcal{N}_I$).}
Remove the element at position $i$ and reinsert it at position $j
\neq i$.  The neighborhood size is $n(n-1)$.  The incremental cost
change (\texttt{deltaInsert}) is computed in $O(n)$.

\subsection{Pivoting Rules}

\paragraph{First Improvement.}
Iterate over neighbors in a fixed order; accept the first move with
$\Delta > 0$ and restart the scan from the beginning.  Terminate when
no improving neighbor is found.

\paragraph{Best Improvement.}
Scan the entire neighborhood; accept the move with the largest positive
$\Delta$.  Terminate when no improving neighbor exists.

Both pivoting rules use function pointers so that any neighborhood can
be passed as a callback without code duplication.

\subsection{Initial Solutions}

\paragraph{Random.}
Generate a uniformly random permutation of $\{1,\ldots,n\}$ using the
Fisher--Yates shuffle and the custom LCG.

\paragraph{Chenery--Watanabe (CW).}
Compute the row sum $a_i = \sum_{j=1}^n c_{ij}$ for each element and
sort elements in decreasing order of $a_i$.  The rationale is that an
element with a large row sum contributes more when placed early in the
permutation, so ranking by $a_i$ yields a greedy approximation to the
optimum.

\subsection{Variable Neighborhood Descent}

Two VND configurations are tested, both initialized with the
Chenery--Watanabe solution and using First Improvement as the internal
pivoting rule:

\begin{itemize}
  \item \textbf{VND-TEI}: neighborhoods cycled in order
        $\mathcal{N}_T \to \mathcal{N}_E \to \mathcal{N}_I$.
  \item \textbf{VND-TIE}: neighborhoods cycled in order
        $\mathcal{N}_T \to \mathcal{N}_I \to \mathcal{N}_E$.
\end{itemize}

The VND procedure optimizes the current solution with respect to the
first neighborhood; upon finding a local optimum, it moves to the next
neighborhood.  If an improvement is found in any non-first neighborhood,
control returns to $\mathcal{N}_T$.  The procedure terminates when the
solution is a local optimum for all three neighborhoods simultaneously.

\subsection{Experimental Setup}

All 14 configurations are run on each of the 78 XLOLIB instances.
Computation time is measured as wall-clock time per run using
\texttt{gettimeofday}.  Statistical comparisons use paired tests (each
algorithm pair is compared instance-by-instance) at significance level
$\alpha = 0.05$: a two-sided paired $t$-test and a Wilcoxon signed-rank
test (with normal approximation).  Both tests agree in all but a handful
of boundary cases, and the Wilcoxon result is reported as primary since
apd values may not be normally distributed.

%─────────────────────────────────────────────────────────────────────────────
\section{Results}
\label{sec:results}

\subsection{Summary Statistics}

Table~\ref{tab:summary} reports, for each of the 14 configurations, the
mean apd, the standard deviation of apd, and the total wall-clock time
summed over all 78 instances.  A lower mean apd means better solution
quality.

\begin{table}[H]
\centering
\caption{Summary of results over 78 XLOLIB instances.
$\bar{\delta}$: mean apd (\%); $\sigma_\delta$: standard deviation of
apd (\%); Time: total time (s) across all instances.}
\label{tab:summary}
\setlength{\tabcolsep}{7pt}
\begin{tabular}{llll r r r}
\toprule
\textbf{Pivoting} & \textbf{Neighborhood} & \textbf{Init}
  & $\bar{\delta}$ & $\sigma_\delta$ & \textbf{Time (s)} \\
\midrule
First  & Transpose & Random  & 34.71 & 5.42  & 0.01 \\
First  & Transpose & Chenery & 16.65 & 2.48  & 0.26 \\
First  & Exchange  & Random  & 34.76 & 5.37  & 0.01 \\
First  & Exchange  & Chenery & 16.76 & 2.49  & 0.26 \\
First  & Insert    & Random  &  5.12 & 0.97  & 243.88 \\
First  & Insert    & Chenery &  5.64 & 1.16  & 144.69 \\
\midrule
Best   & Transpose & Random  & 33.87 & 5.26  & 0.01 \\
Best   & Transpose & Chenery & 15.72 & 2.55  & 0.26 \\
Best   & Exchange  & Random  & 29.06 & 4.50  & 2.60 \\
Best   & Exchange  & Chenery & 14.54 & 2.53  & 1.84 \\
Best   & Insert    & Random  &  5.02 & 1.00  & 26.32 \\
Best   & Insert    & Chenery &  5.85 & 1.16  & 23.61 \\
\midrule
VND-TEI & ---      & Chenery & 16.64 & 2.47  & 0.26 \\
VND-TIE & ---      & Chenery & 16.65 & 2.48  & 0.26 \\
\bottomrule
\end{tabular}
\end{table}

The dominant factor is clearly the \textbf{neighborhood structure}.
Insert-based algorithms achieve mean apd below 6\%, while Transpose and
Exchange algorithms range from 14\% to 35\%.  Within each neighborhood,
the Chenery--Watanabe initialization consistently reduces apd, with the
most pronounced benefit observed for Transpose and Exchange (reductions
of roughly 18--20 percentage points).  For Insert, the initialization
advantage is smaller and partly reversed (random initialization yields
slightly better results than CW for the First Improvement variant).

Best Improvement does not offer consistent gains over First Improvement:
for Insert the differences are marginal, while for Exchange the
Best Improvement variant achieves a modest improvement ($\approx 2\%$
mean apd reduction per pivoting/init combination) at a substantially
higher computational cost.

Both VND variants perform identically to
\textsc{First}+\textsc{Transpose}+\textsc{Chenery} in terms of mean
apd ($\approx 16.6\%$), despite internally exploring Exchange and Insert
neighbors as well.  This indicates that the VND procedure terminates
early at local optima that Insert alone would escape.

\subsection{Statistical Tests}

All 91 pairwise comparisons were performed.
Table~\ref{tab:pairwise} summarizes the most informative comparisons.
Full results are available in the supplementary analysis script.

\begin{table}[H]
\centering
\caption{Selected pairwise Wilcoxon signed-rank tests ($\alpha=0.05$).
$p_W$: Wilcoxon $p$-value; Sig.: statistically significant difference.}
\label{tab:pairwise}
\setlength{\tabcolsep}{6pt}
\begin{tabular}{l l c c}
\toprule
\textbf{Algorithm A} & \textbf{Algorithm B} & $p_W$ & \textbf{Sig.} \\
\midrule
\multicolumn{4}{l}{\textit{Effect of neighborhood (First Improvement, Random init)}} \\
First+Transpose+Rand  & First+Exchange+Rand   & $0.0003$ & Yes \\
First+Transpose+Rand  & First+Insert+Rand     & $<0.0001$ & Yes \\
First+Exchange+Rand   & First+Insert+Rand     & $<0.0001$ & Yes \\
\midrule
\multicolumn{4}{l}{\textit{Effect of initialization (First Improvement)}} \\
First+Transpose+Rand  & First+Transpose+CW    & $<0.0001$ & Yes \\
First+Exchange+Rand   & First+Exchange+CW     & $<0.0001$ & Yes \\
First+Insert+Rand     & First+Insert+CW       & $0.0002$  & Yes \\
\midrule
\multicolumn{4}{l}{\textit{Effect of pivoting rule (Insert neighborhood)}} \\
First+Insert+Rand     & Best+Insert+Rand      & $0.6624$ & No \\
First+Insert+CW       & Best+Insert+CW        & $0.0407$ & Yes \\
\midrule
\multicolumn{4}{l}{\textit{Effect of pivoting rule (Exchange neighborhood)}} \\
First+Exchange+CW     & Best+Exchange+CW      & $<0.0001$ & Yes \\
\midrule
\multicolumn{4}{l}{\textit{VND vs.\ single neighborhood}} \\
First+Transpose+CW    & VND-TEI+CW            & $0.0225$  & Yes \\
First+Transpose+CW    & VND-TIE+CW            & $0.0360$  & Yes \\
VND-TEI+CW            & VND-TIE+CW            & $0.4420$  & No  \\
First+Insert+Rand     & VND-TEI+CW            & $<0.0001$ & Yes \\
\midrule
\multicolumn{4}{l}{\textit{Best insert variants vs.\ VND}} \\
Best+Insert+Rand      & VND-TEI+CW            & $<0.0001$ & Yes \\
Best+Insert+CW        & VND-TEI+CW            & $<0.0001$ & Yes \\
\bottomrule
\end{tabular}
\end{table}

\paragraph{Neighborhood effect.}
All three pairwise comparisons between distinct neighborhoods are
statistically significant ($p < 0.001$ in all cases).  The Insert
neighborhood is the clear winner; the difference between Transpose and
Exchange is smaller but still significant.

\paragraph{Initialization effect.}
The Chenery--Watanabe initialization provides a statistically significant
improvement over random initialization for \emph{all} three
neighborhoods.  The effect is strongest for Transpose and Exchange.
Counterintuitively, for the Insert neighborhood, the random initialization
slightly outperforms CW in the First Improvement setting (mean apd 5.12
vs.\ 5.64), though the difference is statistically significant
($p_W = 0.0002$).  This is likely because the CW heuristic produces a
solution concentrated in a particular region of the search space, from
which Insert search escapes into a different local optimum than the one
reachable from a diverse random start.

\paragraph{Pivoting rule effect.}
For the Insert neighborhood with random initialization, First and Best
Improvement reach solutions of equivalent quality ($p_W = 0.6624$, not
significant), with First Improvement being roughly $10\times$ faster.
For Insert with CW initialization and for Exchange, Best Improvement
achieves slightly but significantly better solutions at a substantial
time cost.  For Transpose, the improvement from Best over First is
negligible.

\paragraph{VND effect.}
Both VND variants significantly outperform
\textsc{First}+\textsc{Transpose}+\textsc{CW} alone ($p_W < 0.05$),
confirming that including the additional neighborhoods is beneficial.
However, VND falls far short of pure Insert-based search (mean apd
$\approx 16.6\%$ vs.\ $\approx 5\%$, $p_W < 0.0001$).  This suggests
that the VND neighborhood cycling terminates at local optima with respect
to all three neighborhoods, yet these solutions are still far from those
found by dedicated Insert search.  The order of neighborhoods (TEI vs.\
TIE) has no statistically significant effect ($p_W = 0.4420$).

\subsection{Computational Cost}

Table~\ref{tab:summary} shows that Insert-based algorithms are by far
the most expensive.  First+Insert+Random requires 244~s total (3.1~s per
instance on average), versus under 0.3~s for all Transpose and Exchange
variants.  Best+Insert is about $10\times$ faster than First+Insert (the
best-improvement scan terminates sooner in practice for this benchmark),
yet achieves comparable or slightly worse solution quality.  VND requires
only 0.26~s total, the same as single Transpose/Exchange with CW, because
it terminates quickly once stuck in a local optimum.

%─────────────────────────────────────────────────────────────────────────────
\section{Conclusion}
\label{sec:conclusion}

This study compared 14 local search configurations for the LOP on the
full XLOLIB benchmark.  The main conclusions are:

\begin{enumerate}
  \item \textbf{Neighborhood choice dominates.}  The Insert neighborhood
        yields mean apd around 5\%, while Transpose and Exchange yield
        15--35\%.  All pairwise differences are statistically significant.
        The LOP search landscape is significantly more amenable to Insert
        moves, as confirmed by the experimental results obtained here.

  \item \textbf{Initialization matters for weak neighborhoods.}
        The Chenery--Watanabe construction heuristic significantly
        reduces apd for Transpose and Exchange ($\approx 18$~pp
        improvement).  For Insert, the benefit is marginal and can
        slightly favor random initialization, suggesting that the
        large Insert neighborhood is less sensitive to starting point.

  \item \textbf{Pivoting rule has modest impact.}
        Best Improvement and First Improvement reach essentially the
        same solution quality for Insert (no significant difference for
        the random-init variant), while Best Improvement incurs
        substantially higher computation time.  For Exchange, Best
        Improvement achieves a small but significant quality gain.

  \item \textbf{VND does not close the gap to Insert search.}
        Despite incorporating Insert moves, VND converges to solutions
        of similar quality to single-neighborhood Transpose search with
        CW initialization ($\approx 16\%$ apd).  This is because VND
        resets to the first neighborhood upon any improvement, causing
        it to terminate at a combined local optimum that differs from
        the pure Insert local optimum.  The neighborhood order (TEI
        vs.\ TIE) has no significant effect.
\end{enumerate}

An effective practical configuration for the LOP under a time budget
is \textbf{First Improvement with the Insert neighborhood and a random
(or CW) initialization}, achieving the lowest mean apd with reasonable
computation time.  Future work could explore iterated local search or
simulated annealing on top of Insert-based local search to further
improve solution quality beyond pure local optima.


\end{document}
